\IfFileExists{stacks-project.cls}{%
\documentclass{stacks-project}
}{%
\documentclass{amsart}
}

% For dealing with references we use the comment environment
\usepackage{verbatim}
\newenvironment{reference}{\comment}{\endcomment}
%\newenvironment{reference}{}{}
\newenvironment{slogan}{\comment}{\endcomment}
\newenvironment{history}{\comment}{\endcomment}

% For commutative diagrams we use Xy-pic
\usepackage[all]{xy}

% We use 2cell for 2-commutative diagrams.
\xyoption{2cell}
\UseAllTwocells

% We use multicol for the list of chapters between chapters
\usepackage{multicol}

% This is generally recommended for better output
\usepackage{lmodern}
\usepackage[T1]{fontenc}

% For cross-file-references
\usepackage{xr-hyper}

% Package for hypertext links:
\usepackage{hyperref}

% For any local file, say "hello.tex" you want to link to please
% use \externaldocument[hello-]{hello}
\externaldocument[introduction-]{introduction}
\externaldocument[conventions-]{conventions}
\externaldocument[sets-]{sets}
\externaldocument[categories-]{categories}
\externaldocument[topology-]{topology}
\externaldocument[sheaves-]{sheaves}
\externaldocument[sites-]{sites}
\externaldocument[stacks-]{stacks}
\externaldocument[fields-]{fields}
\externaldocument[algebra-]{algebra}
\externaldocument[brauer-]{brauer}
\externaldocument[homology-]{homology}
\externaldocument[derived-]{derived}
\externaldocument[simplicial-]{simplicial}
\externaldocument[more-algebra-]{more-algebra}
\externaldocument[smoothing-]{smoothing}
\externaldocument[modules-]{modules}
\externaldocument[sites-modules-]{sites-modules}
\externaldocument[injectives-]{injectives}
\externaldocument[cohomology-]{cohomology}
\externaldocument[sites-cohomology-]{sites-cohomology}
\externaldocument[dga-]{dga}
\externaldocument[dpa-]{dpa}
\externaldocument[sdga-]{sdga}
\externaldocument[hypercovering-]{hypercovering}
\externaldocument[schemes-]{schemes}
\externaldocument[constructions-]{constructions}
\externaldocument[properties-]{properties}
\externaldocument[morphisms-]{morphisms}
\externaldocument[coherent-]{coherent}
\externaldocument[divisors-]{divisors}
\externaldocument[limits-]{limits}
\externaldocument[varieties-]{varieties}
\externaldocument[topologies-]{topologies}
\externaldocument[descent-]{descent}
\externaldocument[perfect-]{perfect}
\externaldocument[more-morphisms-]{more-morphisms}
\externaldocument[flat-]{flat}
\externaldocument[groupoids-]{groupoids}
\externaldocument[more-groupoids-]{more-groupoids}
\externaldocument[etale-]{etale}
\externaldocument[chow-]{chow}
\externaldocument[intersection-]{intersection}
\externaldocument[pic-]{pic}
\externaldocument[weil-]{weil}
\externaldocument[adequate-]{adequate}
\externaldocument[dualizing-]{dualizing}
\externaldocument[duality-]{duality}
\externaldocument[discriminant-]{discriminant}
\externaldocument[derham-]{derham}
\externaldocument[local-cohomology-]{local-cohomology}
\externaldocument[algebraization-]{algebraization}
\externaldocument[curves-]{curves}
\externaldocument[resolve-]{resolve}
\externaldocument[models-]{models}
\externaldocument[functors-]{functors}
\externaldocument[equiv-]{equiv}
\externaldocument[pione-]{pione}
\externaldocument[etale-cohomology-]{etale-cohomology}
\externaldocument[proetale-]{proetale}
\externaldocument[relative-cycles-]{relative-cycles}
\externaldocument[more-etale-]{more-etale}
\externaldocument[trace-]{trace}
\externaldocument[crystalline-]{crystalline}
\externaldocument[spaces-]{spaces}
\externaldocument[spaces-properties-]{spaces-properties}
\externaldocument[spaces-morphisms-]{spaces-morphisms}
\externaldocument[decent-spaces-]{decent-spaces}
\externaldocument[spaces-cohomology-]{spaces-cohomology}
\externaldocument[spaces-limits-]{spaces-limits}
\externaldocument[spaces-divisors-]{spaces-divisors}
\externaldocument[spaces-over-fields-]{spaces-over-fields}
\externaldocument[spaces-topologies-]{spaces-topologies}
\externaldocument[spaces-descent-]{spaces-descent}
\externaldocument[spaces-perfect-]{spaces-perfect}
\externaldocument[spaces-more-morphisms-]{spaces-more-morphisms}
\externaldocument[spaces-flat-]{spaces-flat}
\externaldocument[spaces-groupoids-]{spaces-groupoids}
\externaldocument[spaces-more-groupoids-]{spaces-more-groupoids}
\externaldocument[bootstrap-]{bootstrap}
\externaldocument[spaces-pushouts-]{spaces-pushouts}
\externaldocument[spaces-chow-]{spaces-chow}
\externaldocument[groupoids-quotients-]{groupoids-quotients}
\externaldocument[spaces-more-cohomology-]{spaces-more-cohomology}
\externaldocument[spaces-simplicial-]{spaces-simplicial}
\externaldocument[spaces-duality-]{spaces-duality}
\externaldocument[formal-spaces-]{formal-spaces}
\externaldocument[restricted-]{restricted}
\externaldocument[spaces-resolve-]{spaces-resolve}
\externaldocument[formal-defos-]{formal-defos}
\externaldocument[defos-]{defos}
\externaldocument[cotangent-]{cotangent}
\externaldocument[examples-defos-]{examples-defos}
\externaldocument[algebraic-]{algebraic}
\externaldocument[examples-stacks-]{examples-stacks}
\externaldocument[stacks-sheaves-]{stacks-sheaves}
\externaldocument[criteria-]{criteria}
\externaldocument[artin-]{artin}
\externaldocument[quot-]{quot}
\externaldocument[stacks-properties-]{stacks-properties}
\externaldocument[stacks-morphisms-]{stacks-morphisms}
\externaldocument[stacks-limits-]{stacks-limits}
\externaldocument[stacks-cohomology-]{stacks-cohomology}
\externaldocument[stacks-perfect-]{stacks-perfect}
\externaldocument[stacks-introduction-]{stacks-introduction}
\externaldocument[stacks-more-morphisms-]{stacks-more-morphisms}
\externaldocument[stacks-geometry-]{stacks-geometry}
\externaldocument[moduli-]{moduli}
\externaldocument[moduli-curves-]{moduli-curves}
\externaldocument[examples-]{examples}
\externaldocument[exercises-]{exercises}
\externaldocument[guide-]{guide}
\externaldocument[desirables-]{desirables}
\externaldocument[coding-]{coding}
\externaldocument[obsolete-]{obsolete}
\externaldocument[fdl-]{fdl}
\externaldocument[index-]{index}

% Theorem environments.
%
\theoremstyle{plain}
\newtheorem{theorem}[subsection]{Theorem}
\newtheorem{proposition}[subsection]{Proposition}
\newtheorem{lemma}[subsection]{Lemma}

\theoremstyle{definition}
\newtheorem{definition}[subsection]{Definition}
\newtheorem{example}[subsection]{Example}
\newtheorem{exercise}[subsection]{Exercise}
\newtheorem{situation}[subsection]{Situation}

\theoremstyle{remark}
\newtheorem{remark}[subsection]{Remark}
\newtheorem{remarks}[subsection]{Remarks}

\numberwithin{equation}{subsection}

% Macros
%
\def\lim{\mathop{\mathrm{lim}}\nolimits}
\def\colim{\mathop{\mathrm{colim}}\nolimits}
\def\Spec{\mathop{\mathrm{Spec}}}
\def\Hom{\mathop{\mathrm{Hom}}\nolimits}
\def\Ext{\mathop{\mathrm{Ext}}\nolimits}
\def\SheafHom{\mathop{\mathcal{H}\!\mathit{om}}\nolimits}
\def\SheafExt{\mathop{\mathcal{E}\!\mathit{xt}}\nolimits}
\def\Sch{\mathit{Sch}}
\def\Mor{\mathop{\mathrm{Mor}}\nolimits}
\def\Ob{\mathop{\mathrm{Ob}}\nolimits}
\def\Sh{\mathop{\mathit{Sh}}\nolimits}
\def\NL{\mathop{N\!L}\nolimits}
\def\CH{\mathop{\mathrm{CH}}\nolimits}
\def\proetale{{pro\text{-}\acute{e}tale}}
\def\etale{{\acute{e}tale}}
\def\QCoh{\mathit{QCoh}}
\def\Ker{\mathop{\mathrm{Ker}}}
\def\Im{\mathop{\mathrm{Im}}}
\def\Coker{\mathop{\mathrm{Coker}}}
\def\Coim{\mathop{\mathrm{Coim}}}

% Boxtimes
%
\DeclareMathSymbol{\boxtimes}{\mathbin}{AMSa}{"02}

%
% Macros for moduli stacks/spaces
%
\def\QCohstack{\mathcal{QC}\!\mathit{oh}}
\def\Cohstack{\mathcal{C}\!\mathit{oh}}
\def\Spacesstack{\mathcal{S}\!\mathit{paces}}
\def\Quotfunctor{\mathrm{Quot}}
\def\Hilbfunctor{\mathrm{Hilb}}
\def\Curvesstack{\mathcal{C}\!\mathit{urves}}
\def\Polarizedstack{\mathcal{P}\!\mathit{olarized}}
\def\Complexesstack{\mathcal{C}\!\mathit{omplexes}}
% \Pic is the operator that assigns to X its picard group, usage \Pic(X)
% \Picardstack_{X/B} denotes the Picard stack of X over B
% \Picardfunctor_{X/B} denotes the Picard functor of X over B
\def\Pic{\mathop{\mathrm{Pic}}\nolimits}
\def\Picardstack{\mathcal{P}\!\mathit{ic}}
\def\Picardfunctor{\mathrm{Pic}}
\def\Deformationcategory{\mathcal{D}\!\mathit{ef}}
\usepackage{cleveref}
\crefname{originalremark}{original remark}{original remarks}
\Crefname{originalremark}{Original remark}{Original remarks}
\makeatletter
\newlabel{morphisms-remark-base-change-differential-operators}{{0G46}{source}{}{sourcepointer-1}{}}
\newlabel{morphisms-remark-base-change-differential-operators@cref}{{[originalremark][0][]0G46}{[1][0][]source}{}{}{}}
\makeatother
\crefname{originalsection}{original section}{original sections}
\Crefname{originalsection}{Original section}{Original sections}
\makeatletter
\newlabel{homology-section-truncations}{{0118}{source}{}{sourcepointer-2}{}}
\newlabel{homology-section-truncations@cref}{{[originalsection][0][]0118}{[1][0][]source}{}{}{}}
\makeatother
\crefname{originallemma}{original lemma}{original lemmas}
\Crefname{originallemma}{Original lemma}{Original lemmas}
\makeatletter
\newlabel{lemma-cohomology-de-rham-base-change}{{0FML}{source}{}{sourcepointer-3}{}}
\newlabel{lemma-cohomology-de-rham-base-change@cref}{{[originallemma][0][]0FML}{[1][0][]source}{}{}{}}
\makeatother
\crefname{originallemma}{original lemma}{original lemmas}
\Crefname{originallemma}{Original lemma}{Original lemmas}
\makeatletter
\newlabel{lemma-kunneth-single-sheaf}{{0FU4}{source}{}{sourcepointer-4}{}}
\newlabel{lemma-kunneth-single-sheaf@cref}{{[originallemma][0][]0FU4}{[1][0][]source}{}{}{}}
\makeatother
\crefname{originalequation}{original equation}{original equations}
\Crefname{originalequation}{Original equation}{Original equations}
\makeatletter
\newlabel{equation-de-rham-kunneth}{{0FLR}{source}{}{sourcepointer-5}{}}
\newlabel{equation-de-rham-kunneth@cref}{{[originalequation][0][]0FLR}{[1][0][]source}{}{}{}}
\makeatother
\title{The Kunneth cup product is an isomorphism for bounded complexes flat over an affine base}
\author{The Stacks Project Authors}
\date{stacks-project:perfect.tex:local-snapshot-20261012; GNU FDL 1.2 or later}
\begin{document}
\maketitle
\setcounter{subsection}{1}
\noindent\textit{Editorial note.} All delivered mathematical text is the author's own, copied from the pinned Stacks Project chapter file perfect.tex; no statement or proof step was rewritten, reordered or removed. The counted span is the author environment and its own proof environment. References inside the span are handed to the assembler as navigation stubs that carry the canonical tag of the cited tag only; whether each of them is load-bearing is recorded on the stub itself.\par\medskip
\noindent\textit{Allowed standard prerequisites (editorial).}\begin{itemize}
\item The cup product is functorial, so maps of complexes in the category of differential graded modules give a commutative diagram whose columns are distinguished triangles. used at proof lines 6394-6427. Turns a distinguished triangle of complexes into a relation between the products occurring in the statement.
\item A bounded below complex sits in a distinguished triangle with its truncation above a degree and the degree piece, and the truncation has fewer nonzero terms. invoked at proof lines 6429-6443. Produces the induction that reduces the first complex to a single degree.
\item The statement is of a degreewise nature, so vanishing of the two sides in low degrees lets one replace a bounded below complex by a bounded piece. invoked at proof lines 6445-6490. Produces the same reduction for the second complex, after which only the single degree case is left.
\end{itemize}\medskip
\ifdefined\hypertarget\hypertarget{sourcepointer-1}{}\fi
\noindent\textit{Source reference.} Original remark 0G46: sources/stacks-chow/tags.txt, source lines 1--1. Cited at proof line 6395 for the notation only; every argument carried out in that notation is inside the counted span.\par\medskip
\ifdefined\hypertarget\hypertarget{sourcepointer-2}{}\fi
\noindent\textit{Source reference.} Original section 0118: sources/stacks-chow/tags.txt, source lines 1--1. Cited at proof lines 6434-6436 for the truncation conventions used in the displayed short exact sequence; the triangle and the induction built on it are inside the counted span.\par\medskip
\ifdefined\hypertarget\hypertarget{sourcepointer-3}{}\fi
\noindent\textit{Source reference.} Original lemma 0FML: sources/stacks-perfect/pinned-original-source.tex, source lines 5771--5771. Used at proof lines 6449-6451 to see that the cohomology of the single degree piece has finite tor dimension, which supplies the degree bounds; the degree argument using those bounds is inside the counted span.\par\medskip
\ifdefined\hypertarget\hypertarget{sourcepointer-4}{}\fi
\noindent\textit{Source reference.} Original lemma 0FU4: sources/stacks-perfect/pinned-original-source.tex, source lines 5974--5974. The last sentence at proof line 6495 hands the fully reduced case to that statement; the two reductions that reach it are inside the counted span.\par\medskip
\ifdefined\hypertarget\hypertarget{sourcepointer-5}{}\fi
\noindent\textit{Source reference.} Original equation 0FLR: sources/stacks-perfect/pinned-original-source.tex, source lines 6346--6346. The statement refers to the displayed product by this label; it is a pointer to the diagram, and the proof works with that map directly inside the counted span.\par\medskip
\begin{lemma}
\label{lemma-kunneth-special}
In the situation above the cup product (\ref{equation-de-rham-kunneth})
is an isomorphism in $D(A)$ if the following assumptions hold
\begin{enumerate}
\item $S = \Spec(A)$ is affine,
\item $X$ and $Y$ are quasi-compact with affine diagonal,
\item $\mathcal{F}^\bullet$ is bounded,
\item $\mathcal{G}^\bullet$ is bounded below,
\item $\mathcal{F}^n$ is $S$-flat, and
\item $\mathcal{G}^m$ is $S$-flat.
\end{enumerate}
\end{lemma}

\begin{proof}
We will use the notation $\mathcal{A}_{X/S}$ and $\mathcal{A}_{Y/S}$
introduced in Morphisms, Remark
\ref{morphisms-remark-base-change-differential-operators}.
Suppose that we have maps of complexes
$$
\mathcal{F}_1^\bullet \to
\mathcal{F}_2^\bullet \to
\mathcal{F}_3^\bullet \to
\mathcal{F}_1^\bullet[1]
$$
in the category $\mathcal{A}_{X/S}$.
Then by the functoriality of the cup product
we obtain a commutative diagram
$$
\xymatrix{
R\Gamma(X, \mathcal{F}_1^\bullet)
\otimes_A^\mathbf{L}
R\Gamma(Y, \mathcal{G}^\bullet)
\ar[r] \ar[d] &
R\Gamma(X \times_S Y,
\text{Tot}(\mathcal{F}_1^\bullet \boxtimes \mathcal{G}^\bullet)) \ar[d] \\
R\Gamma(X, \mathcal{F}_2^\bullet)
\otimes_A^\mathbf{L}
R\Gamma(Y, \mathcal{G}^\bullet)
\ar[r] \ar[d] &
R\Gamma(X \times_S Y,
\text{Tot}(\mathcal{F}_2^\bullet \boxtimes \mathcal{G}^\bullet)) \ar[d] \\
R\Gamma(X, \mathcal{F}_3^\bullet)
\otimes_A^\mathbf{L}
R\Gamma(Y, \mathcal{G}^\bullet)
\ar[r] \ar[d] &
R\Gamma(X \times_S Y,
\text{Tot}(\mathcal{F}_3^\bullet \boxtimes \mathcal{G}^\bullet)) \ar[d] \\
R\Gamma(X, \mathcal{F}_1^\bullet[1])
\otimes_A^\mathbf{L}
R\Gamma(Y, \mathcal{G}^\bullet)
\ar[r] &
R\Gamma(X \times_S Y,
\text{Tot}(\mathcal{F}_1^\bullet[1] \boxtimes \mathcal{G}^\bullet))
}
$$
If the original maps form a distinguished triangle in the homotopy category
of $\mathcal{A}_{X/S}$, then
the columns of this diagram form distinguished triangles in $D(A)$.

\medskip\noindent
In the situation of the lemma,
suppose that $\mathcal{F}^n = 0$ for $n < i$. Then we may consider the
termwise split short exact sequence of complexes
$$
0 \to \sigma_{\geq i + 1}\mathcal{F}^\bullet \to
\mathcal{F}^\bullet \to \mathcal{F}^i[-i] \to 0
$$
where the truncation is as in
Homology, Section \ref{homology-section-truncations}.
This produces the distinguished triangle
$$
\sigma_{\geq i + 1}\mathcal{F}^\bullet \to
\mathcal{F}^\bullet \to
\mathcal{F}^i[-i] \to
(\sigma_{\geq i + 1}\mathcal{F}^\bullet)[1]
$$
in the homotopy category of $\mathcal{A}_{X/S}$
where the final arrow is given by the boundary map
$\mathcal{F}^i \to \mathcal{F}^{i + 1}$.
It follows from the discussion above that it suffices to prove the lemma for
$\mathcal{F}^i[-i]$ and $\sigma_{\geq i + 1}\mathcal{F}^\bullet$.
Since $\sigma_{\geq i + 1}\mathcal{F}^\bullet$ has fewer nonzero
terms, by induction, if we can prove the lemma if $\mathcal{F}^\bullet$ is
nonzero only in single degree, then the lemma follows.
Thus we may assume $\mathcal{F}^\bullet$ is nonzero only in one degree.

\medskip\noindent
Assume $\mathcal{F}^\bullet$ is the complex which has an $S$-flat quasi-coherent
$\mathcal{O}_X$-module $\mathcal{F}$ sitting in degree $0$ and is zero in
other degrees. Observe that $R\Gamma(X, \mathcal{F})$ has finite
tor dimension by Lemma \ref{lemma-cohomology-de-rham-base-change} for example.
Say it has tor amplitude in $[i, j]$.
Pick $N \gg 0$ and consider the distinguished triangle
$$
\sigma_{\geq N + 1}\mathcal{G}^\bullet \to
\mathcal{G}^\bullet \to
\sigma_{\leq N}\mathcal{G}^\bullet \to
(\sigma_{\geq N + 1}\mathcal{G}^\bullet)[1]
$$
in the homotopy category of $\mathcal{A}_{Y/S}$. Now observe that both
$$
R\Gamma(X, \mathcal{F})
\otimes_A^\mathbf{L}
R\Gamma(Y, \sigma_{\geq N + 1}\mathcal{G}^\bullet)
\quad\text{and}\quad
R\Gamma(X \times_S Y,
\text{Tot}(\mathcal{F} \boxtimes \sigma_{\geq N + 1}\mathcal{G}^\bullet))
$$
have vanishing cohomology in degrees $\leq N + i$. Thus, using the arguments
given above, if we want to prove our statement in a given degree, then we may
assume $\mathcal{G}^\bullet$ is bounded.
Repeating the arguments above one more time we may also assume
$\mathcal{G}^\bullet$ is nonzero only in one degree.
This case is handled by Lemma \ref{lemma-kunneth-single-sheaf}.
\end{proof}

\end{document}
